\documentclass[11pt]{article} 
\usepackage[french]{babel} 
%FONT
\usepackage[T1]{fontenc} 
\usepackage[utf8]{inputenc}
\usepackage{mathptmx} 
%JUSTIFY 
\usepackage[document]{ragged2e}
%MARGE 
\usepackage{geometry}
\geometry{hmargin=2.5cm,vmargin=2.5cm} 
%Annexe
\usepackage{appendix} 
%Image 
\usepackage{graphicx}
\usepackage[export]{adjustbox}

\begin{document} 
	\begin{titlepage}
		\includegraphics[width=5cm]{img/logo/ecole2.png}.
		%\includegraphics[width=5cm]{img/logo/ecole.png}.
 		\hfill
		\includegraphics[width=5cm]{img/logo/hopital.png}.
		\begin{center} 
			\Huge 
			\vspace{2cm} 
			Rapport de projet\\
			\vspace{1cm} 
			\normalsize ISIS - FIA 3 - Promotion 2022\\
			\vspace{3cm} 
			Hopital de Narbonne 
		\end{center} 
		\vfill
		\begin{tabular}{p{6cm}p{6cm}p{6cm}} 
			\underline{Apprenti} & \underline{Maitre apprentissage} & \underline{Tuteur école}\\ 
			& & \\ 
			Martin BORDE & Hélène LHERBETTE & Laurent GRÉGOIRE \\
		\end{tabular}
	\end{titlepage} 
%\title{Rapport de projet\\ 
%MoléculesOnéreuses} 
%\author{Martin BORDE\\ 
% Ecole d'ingénieur ISIS,\\ 
%Castres,\\ 
% \texttt{}} % % 
%\maketitle \newpage 

%Remerciement
Remerciement \newpage

 %table des matières 
\tableofcontents
\newpage 

\justify

%Corps de rapport

\section{Un rapport d'action}
%DEMANDE DU SUJET
%Il s'agit d’un résumé en 4 pages de l'ensemble du projet,
%destiné à un décideur. Il doit présenter de façon claire et
%synthétique le problème posé, les objectifs fixés, les résultats
%obtenus et les principaux éléments de management du projet. Il
%faut imaginer que ce rapport d'action sera le seul document lu
%par le directeur de l'entreprise de l'apprenti pour juger de son
%travail. La solution proposée et/ou mise en œuvre ne doit pas
%être développée dans le détail.

	\subsection{Problème}
	\subsection{Objectifs fixés}
	\subsection{Résultats}
	\subsection{Management du projet}

\section{Un rapport d'étude}
	\subsection{Base de données}

%DEMANDE DU SUJET
%Le rapport d'étude (20 à 30 pages), destiné à la maîtrise
%d'œuvre, doit faire apparaître tous les détails utiles à la
%compréhension du problème posé jusqu'à sa résolution. La
%solution proposée et/ou mise en œuvre doit être développée dans
%le détail pour les aspects techniques.

\section{Une description du management du projet}
	\subsection{GANTT}
	\subsection{Réunions}
%DEMANDE DU SUJET
%Dans cette partie (4 à 6 pages) doivent figurer tous les
%éléments de management du projet : schéma heuristique,
%organigramme des tâches, planning prévisionnel et réalisé avec
%analyse des écarts, démarche retenue, outils utilisés, relations
%avec le maître d'apprentissage et l’équipe projet, ....

\section{Un bilan personnel du projet professionnel}

	\subsection{Mise en pratique}
%DEMANDE DU SUJET
%Il est demandé à l'apprenti dans cette partie de construire son
%projet professionnel à partir de l'expérience de l'année passée
%dans l'entreprise (en 2 à 6 pages). Par ailleurs, l'apprenti
%doit rédiger la liste des compétences qu'il a développées dans
%l'entreprise en s'appuyant sur des référentiels5


\begin{appendix} \end{appendix}

\end{document}

